\documentclass[12pt,a4paper]{article}
\usepackage{ucebnice}
\usepackage{polynom}

\title{Kvadratické rovnice a nerovnice}
\author{Ondřej Škubala}
\date{\today}

\begin{document}

\maketitle
\newpage

\section{Úvod}

\begin{quotebox}
{„Žádná rovnice není nemožná, dokud neztrácíme trpělivost hledat její řešení.“}
{al-Chvárizmí (9. století)}
\end{quotebox}

\begin{historicalbox}
\textbf{Od geometrie k algebře}

Počátky kvadratických rovnic sahají více než tři tisíce let do minulosti. 
V době, kdy ještě neexistoval symbolický jazyk matematiky, 
zapisovali staří Babyloňané své úvahy na hliněné tabulky klínovým písmem a 
řešili úlohy, které bychom dnes popsali rovnicemi typu 
\(x^2 + bx = c\) či \(x^2 = bx + c\). V jedné z dochovaných tabulek, označované 
jako \textit{YBC 6967}, je popsán postup, jak najít délku strany čtverce, 
pokud známe jeho plochu a přebytek určité části – řešení vychází z doplnění čtverce, 
tedy stejné myšlenky, na které je založen i dnešní vzorec pro kvadratickou rovnici.  
Babyloňané ovšem neznali symboly ani záporná čísla, 
a tak vše popisovali slovně a geometricky: 
\textit{„vezmi polovinu počtu, násob sama sebou, přidej a odečti, dokud nedostaneš čtverec“.} 
Přesto byl výsledek překvapivě přesný a správný.

\medskip
Ve starověkém Egyptě se podobné úlohy objevují v tzv. 
Rhindově papyru (asi 1650 př. n. l.), kde jsou rovnice vyjádřeny prostřednictvím 
„falešného předpokladu“ – Egypťané zvolili odhad hodnoty, spočítali výsledek a 
poté upravili chybu poměrem. I když nešlo o algebraické myšlení v dnešním smyslu, 
jednalo se o první systematické metody řešení rovnic.

\medskip
\textbf{Řecko – geometrické období}

Antičtí Řekové přenesli problém kvadratických vztahů do oblasti geometrie. 
Eukleidés ve svých \textit{Základech} dokázal řadu vět o rovnosti ploch, 
které dnes překládáme jako algebraické identity, například \((a+b)^2 = a^2 + 2ab + b^2\). 
Herón Alexandrijský i Diofantos z Alexandrie se již přiblížili k řešení rovnic 
se symbolickými náznaky, přestože symbol „x“ ještě neexistoval. 
Každé řešení bylo v jejich pojetí konstrukcí s pravítkem a kružítkem – 
nalezení délky úsečky, která vyhovuje danému vztahu.

\medskip
\textbf{Arabská matematika – zrod algebry}

K zásadnímu posunu došlo v 9. století v Bagdádu, 
který se tehdy stal centrem vzdělanosti islámského světa. 
Právě zde působil učenec \textbf{Muhammad ibn Músá al-Chvárizmí}, 
autor knihy \textit{Al-jabr wal-muqabala}. 
Slovo \textit{al-jabr} znamená „doplnění“ (např. chybějící části čtverce) 
a \textit{al-muqabala} „vyrovnání“. V této práci popsal systematický postup, 
jak řešit šest základních typů kvadratických rovnic, které tehdejší svět znal.  
Jeho přístup byl čistě slovní, avšak zcela obecný – al-Chvárizmí dokázal, 
že každou kvadratickou rovnici lze převést na standardní tvar a 
řešit doplněním na čtverec. 
A právě z názvu \emph{al-jabr} se později vyvinulo naše slovo \textbf{algebra}.

\textbf{Al-Chvárizmí: slovní řešení kvadratické úlohy (9. století)}

„\textit{První z úloh. Řeknou-li ti: Rozdělil jsi deset na dvě části, vynásobil jsi jednu z částí druhou a dále jsi vynásobil jednu z nich samu sebou, potom tento součin sama se sebou se stal rovným čtyřnásobku součinu obou částí.}

\textit{Pravidlo je následující: Vezmi jednu z částí jako věc, potom druhá je deset bez věci. Vynásob věc krát deset bez věci, obdržíš deset věcí bez kvadrátu. Dále to násob čtyřmi, jak ti bylo řečeno: čtyřikrát. Obdržíš čtyřicet věcí bez čtyř kvadrátů.}

\textit{Potom násob jednu část druhou, to jest čtyřnásobek součinu obou částí. Potom násob věc krát věc, to jest věc samu sebou. Obdržíš kvadrát. Doplň k tomu čtyřikrát věcí obdržíš: čtyřicet věcí bez čtyř kvadrátů a přičti ke kvadrátu. Obdržíš: čtyřicet věcí bez tří kvadrátů.}

\textit{Proto je jeden kvadrát roven čtyřiceti věcem bez čtyř kvadrátů odtěžených. Čtyřikrát věcí rovno pěti kvadrátům.}“

\smallskip
\textit{Zdroj:}  
Al-Chvárizmí: \textit{Algebraický a aritmetický traktát}.  
Český překlad, Praha 2023.  
ISBN: 978-980-87269-01-5, str. 134.


\textbf{Stejná úloha řešená „al-Chvárizmího metodou“}

Uvažujeme stejnou situaci jako předtím: číslo \(10\) dělíme na dvě části \(x\) a \(10-x\) tak,  
že součin jedné části se sebou samou je roven čtyřnásobku součinu obou částí:
\[
x^2 = 4x(10-x).
\]

Nejprve upravíme rovnici tak, aby měla tvar „kvadrát rovný věcem“, jak to dělá al-Chvárizmí.  
Roznásobíme pravou stranu:
\[
x^2 = 40x - 4x^2.
\]

Převedeme všechny členy na jednu stranu:
\[
5x^2 - 40x = 0.
\]

Vydělíme pěti:
\[
x^2 - 8x = 0.
\]

Al-Chvárizmí by tuto rovnici popsal slovy jako „kvadrát rovný osmi věcem“ –  
kvadrát \((x^2)\) a osm násobků věci \((8x)\).

Dále použije pravidlo doplnění na čtverec:

„Vezmi polovinu počtu věcí, vynásob ji sebou samou a přičti tuto hodnotu na obě strany.“

Polovina z \(8\) je \(4\), její čtverec je
\[
4^2 = 16.
\]

Přičteme \(16\) na obě strany:
\[
x^2 - 8x + 16 = 16.
\]

Levá strana je nyní dokonalý čtverec:
\[
(x - 4)^2 = 16.
\]

Odmocníme obě strany:
\[
\sqrt{(x - 4)^2} = \sqrt{16},
\]
\[
|x - 4| = 4.
\]

V klasickém al-Chvárizmího pojetí se bere kladná varianta, takže
\[
x - 4 = 4,
\]
\[
x = 8.
\]

Druhá část je pak
\[
10 - x = 2.
\]

Tím jsme stejnou úlohu vyřešili metodou, kterou al-Chvárizmí popisuje slovně:  
převedením rovnice na tvar „kvadrát rovný věcem“, doplněním na čtverec a odmocněním.



\medskip
Arabská matematika navíc položila základy symboliky a 
přenesla do Evropy nejen řecké spisy, ale i vlastní poznatky o číslech, úhlech a rovnicích. 
Přes Španělsko a severní Afriku se tyto ideje dostaly do Itálie, 
kde ve 13.–15. století vznikají první učebnice obchodní matematiky. 
Mezi nimi vynikl \textbf{Leonardo z Pisy}, známý jako Fibonacci, 
který do Evropy přinesl indicko-arabské číslice i základní principy řešení rovnic.

\medskip
\textbf{Renesance a novověk – přechod k symbolické algebře}

V 16. století se v Itálii objevují první pokusy o řešení rovnic třetího a 
čtvrtého stupně (Cardano, Tartaglia, Ferrari - ne ten Ferrari s traktory), 
avšak kvadratická rovnice již byla plně zvládnutá disciplína. 
K zásadní proměně dochází s dílem \textbf{Françoise Vièta} a 
později \textbf{Reného Descarta}, kteří zavedli symbolický jazyk, 
jak jej používáme dodnes: písmena \(a, b, c\) označují koeficienty, 
\(x\) neznámou a rovnice se stává čistě algebraickým objektem, 
nezávislým na geometrické představě.

\medskip
V 17. století se díky rozvoji analytické geometrie a diferenciálního počtu
 začala kvadratická rovnice chápat i jako nástroj k popisu fyzikálních jevů. 
 Parabola, jejíž rovnice má tvar \(y = ax^2 + bx + c\), 
 se objevuje ve studiích \textbf{Galilea Galileiho} o pohybu těles, 
 v optice i balistice. Od té doby se kvadratické vztahy staly neoddělitelnou součástí přírodních věd.


\textbf{Zajímavosti a souvislosti}

Paraboly nalezneme nejen v matematice, ale i v architektuře a umění. Klenby mostů, tvary reflektorů, dráhy planet či dokonce křivky při hodu míčem – to vše lze popsat právě kvadratickými rovnicemi. Moderní architekt Antonio Gaudí využíval parabolické tvary při návrhu svých staveb v Barceloně, protože rozkládají síly přirozeným způsobem. V hudbě se kvadratické vztahy objevují ve formě změn amplitudy a frekvence, ve fyzice pak popisují potenciální energii pružiny \(E = \tfrac{1}{2}kx^2\).


Kvadratická rovnice se tak stala nejen matematickým objektem, ale i symbolem obecného principu harmonie mezi částmi a celkem – mezi přímkou a křivkou, rovnováhou a extrémem, mezi jednoduchostí zápisu a hloubkou významu, který dokáže vyjádřit.

\end{historicalbox}



\section{Kvadratická rovnice}
Navazujíc na historické kořeny a význam kvadratických vztahů se nyní zaměříme na 
samotnou podstatu \textit{kvadratických rovnic} a \textit{nerovnic}.  
Zatímco v předchozích kapitolách jsme zkoumali kvadratickou funkci především 
z grafického hlediska a sledovali, jak se mění její tvar v závislosti na parametrech, 
nyní se zaměříme na algebraickou stránku věci — na hledání takových hodnot proměnné, 
které dané rovnici vyhovují.

Kvadratické rovnice představují základní typ nelineárních rovnic a 
objevují se ve všech oblastech matematiky i fyziky, 
od výpočtu dráhy tělesa (pokud jste hráli Kerbal Space Program, tak tušítes) po optimalizační úlohy v ekonomii. 
Jejich řešení vede přirozeně k myšlence doplnění na čtverec a k zavedení tzv. 
\textbf{diskriminantu}, který umožňuje rozhodnout, kolik \textbf{reálných řešení} rovnice má.  
Tato kapitola tak tvoří přirozený most mezi světem funkcí a světem rovnic — 
mezi geometrickým a algebraickým pohledem na tentýž problém.

\begin{definitionbox}
\textbf{Kvadratická rovnice} je každá rovnice, kterou lze po úpravách zapsat ve tvaru
\[
ax^2 + bx + c = 0,
\]
kde \(a, b, c \in \mathbb{R}\) a zároveň \(a \neq 0\).

Čísla \(a, b, c\) nazýváme \textbf{koeficienty kvadratické rovnice}, přičemž:
\begin{itemize}
  \item \(a\) je \textbf{kvadratický koeficient},
  \item \(b\) je \textbf{lineární koeficient},
  \item \(c\) je \textbf{absolutní člen}.
\end{itemize}

Hledáme všechna reálná čísla \(x\), pro která je daná rovnice splněna.  
Tato čísla nazýváme \textbf{kořeny rovnice} a jejich množinu značíme \(K\).
\end{definitionbox}


Po zavedení obecného tvaru kvadratické rovnice si nyní všimněme, 
že podle přítomnosti či nepřítomnosti jednotlivých členů 
rozlišujeme tři základní typy kvadratických rovnic:

\begin{itemize}
  \item \textbf{ryze kvadratické rovnice} – tvaru \(ax^2 + c = 0\),
  \item \textbf{neúplné kvadratické rovnice} – tvaru \(ax^2 + bx = 0\),
  \item \textbf{úplné kvadratické rovnice} – tvaru \(ax^2 + bx + c = 0\).
\end{itemize}


\subsubsection*{Ryze kvadratická rovnice}
Každý z těchto typů má svůj specifický postup řešení. 
Začněme nejjednodušším případem – \textbf{ryze kvadratickou rovnicí}.



Rovnici
\[
ax^2 + c = 0
\]
můžeme upravit následovně:
\[
ax^2 = -c \quad \Rightarrow \quad x^2 = -\frac{c}{a}.
\]
Označme výraz \(-\tfrac{c}{a}\) jako \(k\). Potom dostáváme jednodušší tvar
\[
x^2 = k,
\]
který lze přepsat i ve tvaru
\[
x^2 - k = 0.
\]

Z algebry víme, že pro každý výraz tvaru \(a^2 - b^2\) platí vzorec
\[
a^2 - b^2 = (a - b)(a + b),
\]
a tudíž můžeme rozepsat náš výraz jako:
\[
x^2 - k = 0.
\]
\[
(x - \sqrt{k})(x + \sqrt{k}) = 0.
\]
Aby byl součin dvou činitelů nulový, musí být alespoň jeden z nich roven nule, 
proto tedy platí:
\[
x_1 = \sqrt{k}, \qquad x_2 = -\sqrt{k}.
\]



Tentýž výsledek lze zapsat i stručnějším způsobem.  
Protože u kvadratické funkce platí, že pro opačné hodnoty proměnné \(x\) 
získáváme stejné hodnoty výrazu \(x^2\), 
můžeme psát:
\[
x^2 = k \quad \Rightarrow \quad |x| = \sqrt{k}.
\]
A tedy:
\[
x_1 = \sqrt{k}, \qquad x_2 = -\sqrt{k}.
\]
Což po zpětném dosazením z $k=-\frac{c}{a}$ vypadá:

\[
x_1 = \sqrt{-\frac{c}{a}}, \qquad x_2 = -\sqrt{-\frac{c}{a}}.
\]



Pokud by však hodnota \(k\) byla záporná (tj. \(-\tfrac{c}{a} < 0\)), 
rovnice by neměla žádné reálné řešení, 
neboť druhá mocnina reálného čísla nemůže být záporná. 
V takovém případě bychom řešení hledali v oboru komplexních čísel (s čímž se ještě setkáme v dalším ročníku).

\begin{examplebox}
\textbf{Ryze kvadratická rovnice}

Vyřešte rovnic:
\[
25x^2 + 1 = 0
\]
\[
x^2 = -\frac{1}{25}.
\]
\[
x = \sqrt{-\frac{1}{25}}.
\]
Protože druhá mocnina reálného čísla nemůže být záporná, rovnice nemá reálné řešení.
\[
K = \emptyset.
\]
\end{examplebox}

Když nyní tuto rovnici pozměníme a to tak, že tentokrát kvadratický a absolutní člen budeme od sebe odečíta. A ukážeme si, jak se tato změna projeví na výsledku kořenu rovnice.

\begin{examplebox}
Vyřešte rovnici:
\[
25x^2 - 1 = 0
\] 
\[
x^2 = \frac{1}{25}.
\]
Odmocníme obě strany:
\[
\sqrt{x^2} = \sqrt{\frac{1}{25}}.
\]
\[
|x| = \frac{1}{5}.
\]
\textbf{Množina řešení:}
\[
K = \left\{ -\tfrac{1}{5}, \tfrac{1}{5} \right\}.
\]
\end{examplebox}




\subsubsection*{Neúplná kvadratická rovnice}
Nyní se podívejme na druhý typ, tzv. \textbf{neúplnou kvadratickou rovnici}, 
která má tvar
\[
ax^2 + bx = 0.
\]

Všimněme si, že obě části rovnice obsahují proměnnou \(x\).  
Proto můžeme společného činitele \(x\) vytknout:
\[
x(ax + b) = 0.
\]

Tím jsme rovnici rozložili na součin dvou výrazů.  
Aby byl tento součin nulový, musí být alespoň jeden z činitelů roven nule, 
tedy platí:
\[
x = 0 \quad \text{nebo} \quad ax + b = 0.
\]

Z druhé rovnice \(ax + b = 0\) pak snadno dostáváme
\[
x = -\frac{b}{a}.
\]



Tím jsme získali dva kořeny:
\[
x_1 = 0, \quad x_2 = -\frac{b}{a}.
\]



Pokud by však platilo \(b = 0\), potom by se rovnice zjednodušila na \(ax^2 = 0\), 
z čehož plyne \(x = 0\).  
V tomto případě má rovnice jediný kořen, který je však tzv. 
\textbf{dvojnásobným kořenem}, 
protože odpovídá situaci, kdy se parabola dotýká osy \(x\) právě v jednom bodě.



Shrneme-li, neúplné kvadratické rovnice se řeší snadno rozkladem na součin, 
přičemž vždy jeden z kořenů bývá roven nule. 
Tento tvar se velmi často objevuje i v praxi, například při výpočtech dráhy nebo síly, 
kde je jeden z faktorů společný pro více členů rovnice.

\begin{examplebox}
\textbf{Neúplná kvadratická rovnice}

Vyřešte rovnici:
\[
2x^2 + \tfrac{3}{2}x = 0.
\]

Oba členy obsahují proměnnou \(x\), proto ji vytkneme:
\[
x\left(2x + \tfrac{3}{2}\right) = 0.
\]

Aby byl součin nulový, musí být alespoň jeden z činitelů nulový:
\[
x = 0 \quad \text{nebo} \quad 2x + \tfrac{3}{2} = 0.
\]

Z druhé rovnice získáme:
\[
2x = -\tfrac{3}{2}.
\]
\[
x = -\tfrac{3}{4}.
\]

\textbf{Množina řešení:}
\[
K = \left\{ 0,\; -\tfrac{3}{4} \right\}.
\]
\end{examplebox}





\subsubsection*{Úplná kvadratická rovnice}

Třetím a nejdůležitějším typem je tzv. \textbf{úplná kvadratická rovnice}, 
která obsahuje všechny tři členy:
\[
ax^2 + bx + c = 0,
\]
kde \(a, b, c \in \mathbb{R}\) a \(a \neq 0\).

Tento typ rovnice nelze jednoduše rozložit na součin, 
protože se zde současně objevují členy \(x^2\), \(x\) i konstantní člen.  
Abychom ji dokázali vyřešit, použijeme metodu, kterou už ve svém díle 
\textit{Al-jabr wal-muqabala} popisoval al-Chvárizmí — tzv. \textbf{doplnění na čtverec}.


\textit{Hic sunt leones et dracones:}

Začněme s obecnou rovnicí
\[
ax^2 + bx + c = 0.
\]
Nejprve ji vydělíme koeficientem \(a\), aby byl člen u \(x^2\) rovný jedné:
\[
x^2 + \frac{b}{a}x + \frac{c}{a} = 0.
\]
Nyní si z prvních dvou členů všimněme, že by mohly tvořit druhou mocninu nějakého výrazu.
Připomeňme si vzorec:
\[
(x + p)^2 = x^2 + 2px + p^2.
\]
Abychom mohli použít tento tvar, potřebujeme, aby platilo \(2p = \tfrac{b}{a}\).  
Odtud vyplývá \(p = \tfrac{b}{2a}\).



K oběma stranám rovnice tedy přičteme výraz 
\(\left(\tfrac{b}{2a}\right)^2\), abychom mohli zleva vytvořit dokonalý čtverec:
\[
x^2 + \frac{b}{a}x + \frac{c}{a} + \left(\frac{b}{2a}\right)^2 = \left(\frac{b}{2a}\right)^2.
\]
\[
x^2 + \frac{b}{a}x + \left(\frac{b}{2a}\right)^2 = \left(\frac{b}{2a}\right)^2 - \frac{c}{a}.
\]
Levá strana je nyní dokonalý čtverec:
\[
\left(x + \frac{b}{2a}\right)^2 = \left(\frac{b}{2a}\right)^2 - \frac{c}{a}.
\]
A ještě upravíme pravou stranu rovnice:
\[
\left(x + \frac{b}{2a}\right)^2 = \frac{b^2 - 4ac}{4a^2}.
\]

Odmocníme obě strany:
\[
\sqrt{\left(x + \frac{b}{2a}\right)^2} = \sqrt{\frac{b^2 - 4ac}{4a^2}}.
\]
\[
\bigg|x + \frac{b}{2a}\bigg| = \frac{\sqrt{b^2 - 4ac}}{2a}.
\]
\[
x + \frac{b}{2a} = \pm \frac{\sqrt{b^2 - 4ac}}{2a}.
\]
A po odečtení \(\tfrac{b}{2a}\) dostáváme:
\[
x  = - \frac{b}{2a} \pm \frac{\sqrt{b^2 - 4ac}}{2a},
\]
což lze upravit na tvar:
\[
x = \frac{-b \pm \sqrt{b^2 - 4ac}}{2a}.
\]



Tento výsledek se nazývá \textbf{vzorcem pro kořeny kvadratické rovnice}.  
Výraz pod odmocninou
\[
D = b^2 - 4ac
\]
nazýváme \textbf{diskriminant} a jeho hodnota rozhoduje o počtu reálných řešení:
\[
\begin{array}{rcl}
D > 0 & \Rightarrow & \text{dva různé reálné kořeny},\\[4pt]
D = 0 & \Rightarrow & \text{jeden dvojnásobný kořen},\\[4pt]
D < 0 & \Rightarrow & \text{žádné reálné řešení (komplexní)}.
\end{array}
\]


\begin{definitionbox}
\textbf{Vzor pro kořeny kvadratické rovnice:}

\medskip
\[
x_{1,2} = \frac{-b \pm \sqrt{D}}{2a},  \text{ kde } D = b^2 - 4ac
\]

\bigskip
Výraz pod odmocninou
\[
D = b^2 - 4ac
\]
nazýváme \textbf{diskriminant} a jeho hodnota rozhoduje o počtu reálných řešení:
\[
\begin{array}{rcl}
D > 0 & \Rightarrow & \text{dva různé reálné kořeny},\\[4pt]
D = 0 & \Rightarrow & \text{jeden dvojnásobný kořen},\\[4pt]
D < 0 & \Rightarrow & \text{žádné reálné řešení (pouze komplexní)}.
\end{array}
\]
\end{definitionbox}


Odvozením tohoto vzorce jsme završili řešení všech typů kvadratických rovnic. 


\begin{examplebox}
\textbf{Úplná kvadratická rovnice}

Vyřešte rovnici:
\[
x^2 + 2x - 8 = 0.
\]

Určíme koeficienty: \(a = 1,\; b = 2,\; c = -8.\)


Spočítáme diskriminant:
\[
D = b^2 - 4ac = 2^2 - 4 \cdot 1 \cdot (-8) = 4 + 32 = 36.
\]

\[
x_{1,2} = \frac{-b \pm \sqrt{D}}{2a} = \frac{-2 \pm 6}{2}.
\]

\[
x_1 = 2, \quad x_2 = -4.
\]

\textbf{Množina řešení:}
\[
K = \{-4,\; 2\}.
\]
\end{examplebox}



\begin{examplebox}
\textbf{Úplná kvadratická rovnice bez reálného řešení}

Vyřešte rovnici:
\[
2x^2 - 4x + 3 = 0.
\]

Určíme koeficienty: \(a = 2,\; b = -4,\; c = 3.\)


Spočítáme diskriminant:
\[
D = b^2 - 4ac = (-4)^2 - 4 \cdot 2 \cdot 3 = 16 - 24 = -8.
\]

Protože \(D < 0\), rovnice nemá žádné reálné řešení.

\textbf{Množina řešení:}
\[
K = \emptyset.
\]
\end{examplebox}


\newpage
\subsection{Vztahy mezi kořeny a koeficienty kvadratické rovnice}

Uvažujme kvadratickou rovnici
\[
ax^2 + bx + c = 0,
\]
kde \(a, b, c \in \mathbb{R}\), \(a \neq 0\) a předpokládejme, že má dva reálné kořeny \(x_1, x_2\).  
Z obecného vzorce pro kořeny víme, že:
\[
x_1 = \frac{-b + \sqrt{b^2 - 4ac}}{2a}, \qquad x_2 = \frac{-b - \sqrt{b^2 - 4ac}}{2a}.
\]


\textbf{Součet kořenů:}

\[
x_1 + x_2 = \frac{-b + \sqrt{b^2 - 4ac}}{2a} + \frac{-b - \sqrt{b^2 - 4ac}}{2a}.
\]
Sečtením čitatelů dostaneme:
\[
x_1 + x_2 = \frac{-b + \sqrt{b^2 - 4ac} - b - \sqrt{b^2 - 4ac}}{2a}.
\]
Odmocniny se odečtou:
\[
x_1 + x_2 = \frac{-2b}{2a}.
\]
Zjednodušením obdržíme:
\[
x_1 + x_2 = -\frac{b}{a}.
\]


\textbf{Součin kořenů:}

\[
x_1 \cdot x_2 = \frac{\left(-b + \sqrt{b^2 - 4ac}\right)\left(-b - \sqrt{b^2 - 4ac}\right)}{4a^2}.
\]
Využijeme vzorec pro součin dvojčlenů \( (u+v)(u-v) = u^2 - v^2 \):
\[
x_1 \cdot x_2 = \frac{(-b)^2 - (\sqrt{b^2 - 4ac})^2}{4a^2}.
\]
Po úpravě:
\[
x_1 \cdot x_2 = \frac{b^2 - (b^2 - 4ac)}{4a^2} = \frac{4ac}{4a^2} = \frac{c}{a}.
\]

Tyto vztahy tedy platí pro každou kvadratickou rovnici s reálnými kořeny:
\[
\begin{aligned}
x_1 + x_2 &= -\frac{b}{a}, \\[6pt]
x_1 \cdot x_2 &= \frac{c}{a}.
\end{aligned}
\]


\subsubsection*{Viètovy vzroce}
Tyto elegantní vztahy mezi kořeny a koeficienty kvadratické rovnice poprvé formuloval na konci 16. století 
\textbf{François Viète} (1540 - 1603), významný francouzský matematik, 
který zavedl systematické používání písmen pro neznámé i známé veličiny.  
Díky němu se algebra stala obecnou symbolickou řečí matematiky a právě tyto vztahy —
 známé jako \textbf{Viètovy vztahy} — umožňují nejen ověřit správnost výpočtu, 
 ale i sestavit kvadratickou rovnici, známe-li její kořeny.

\begin{definitionbox}
\textbf{Viètovy vztahy:}
\[
x_1 + x_2 = -\frac{b}{a}, \qquad x_1 \cdot x_2 = \frac{c}{a}.
\]
\end{definitionbox}

\medskip
Tyto vztahy lze využít i opačně:  
Pokud známe kořeny \(x_1, x_2\), můžeme rovnici zpětně sestavit ve tvaru
\[
a(x - x_1)(x - x_2) = 0,
\]
který po roznásobení dává
\[
ax^2 - a(x_1 + x_2)x + a x_1 x_2 = 0.
\]
Porovnáním s obecným tvarem \(ax^2 + bx + c = 0\) znovu vidíme, že
\[
-b = a(x_1 + x_2) \quad \text{a} \quad c = a x_1 x_2.
\]

\begin{examplebox}
\textbf{Ukázkový příklad: Sestavení kvadratické rovnice z daných kořenů}

Najděme všechny kvadratické rovnice, jejichž kořeny jsou 
\[
x_1 = -1, \qquad x_2 = \tfrac{3}{4}.
\]

Každá kvadratická rovnice s kořeny \(x_1\) a \(x_2\) musí mít tvar
\[
a(x - x_1)(x - x_2) = 0,
\]
kde \(a \in \mathbb{R}\) a \(a \neq 0\).  
Tento zápis zachycuje celou rodinu rovnic, které mají stejné kořeny, protože volba koeficientu \(a\) odpovídá jen násobení rovnice nenulovou konstantou.

Po dosazení známých kořenů získáme
\[
a(x + 1)\left(x - \tfrac{3}{4}\right) = 0.
\]

Roznásobením dostáváme
\[
a\left(x^2 + x - \tfrac{3}{4}x - \tfrac{3}{4}\right)
=
a\left(x^2 + \tfrac{1}{4}x - \tfrac{3}{4}\right)
= 0.
\]

Všechny kvadratické rovnice s těmito kořeny tedy mají tvar
\[
ax^2 + \tfrac{a}{4}x - \tfrac{3a}{4} = 0,
\]
kde \(a\neq 0\).

Pro pohodlnější zápis můžeme rovnici vynásobit čtyřmi a dostaneme ekvivalentní rodinu rovnic
\[
4ax^2 + ax - 3a = 0.
\]

Zvolíme-li například \(a = 4\), dostáváme rovnici bez zlomků
\[
4x^2 + x - 3 = 0,
\]
která má kořeny přesně \(-1\) a \(\tfrac{3}{4}\), stejně jako všechny ostatní rovnice této rodiny.

\end{examplebox}




\section{Kvadratická rovnice s parametrem}

V dosud řešených úlohách byly koeficienty kvadratické rovnice pevně danými reálnými čísly, 
zatímco v této části budeme uvažovat obecnější situaci, kdy se některý z koeficientů 
stává proměnným \textit{parametrem}.  
Taková rovnice nemá jedno konkrétní řešení, ale celou rodinu řešení, 
protože její tvar i počet kořenů závisí na hodnotě parametru.  
Cílem nebude vždy nalézt konkrétní kořeny, nýbrž určit, 
\textit{pro které hodnoty parametru má rovnice řešení, kolik jich je, případně jaké vlastnosti mají}.  

Kvadratické rovnice s parametrem tak představují přirozený most mezi obyčejným počítáním a 
obecnějším algebraickým myšlením, v němž nepracujeme s jedním konkrétním případem, 
ale s celou množinou případů současně.

\begin{definitionbox}
\textbf{Kvadratická rovnice s parametrem} je rovnice tvaru
\[
a(p)x^2 + b(p)x + c(p) = 0,
\]
kde alespoň jeden z koeficientů závisí na parametru \(p \in \mathbb{R}\), 
přičemž stále platí podmínka \(a(p) \neq 0\).
\end{definitionbox}

Výhodou parametrického zápisu je možnost zkoumat, 
jak se mění počet řešení rovnice při různých hodnotách parametru.  
Nápovědí je opět diskriminant:
\[
D(p) = b(p)^2 - 4a(p)c(p),
\]
který sám o sobě bývá funkcí parametru.  
Počet kořenů pak určujeme stejně jako u běžné kvadratické rovnice:

\[
\begin{array}{lll}
D(p) > 0 & \Rightarrow & \text{rovnice má dva různé reálné kořeny},\\[4pt]
D(p) = 0 & \Rightarrow & \text{rovnice má jeden (dvojnásobný) kořen},\\[4pt]
D(p) < 0 & \Rightarrow & \text{rovnice nemá reálné kořeny}.\\
\end{array}
\]

To však nebývá vše — často jsme po studentovi požádáni, 
aby našel takové hodnoty parametru, 
pro něž má rovnice např. \textit{kladné řešení}, \textit{záporné řešení}, 
\textit{kořen v určitém intervalu}, 
či aby měla \textit{právě jeden kořen}, popřípadě aby byla \textit{řešitelná}.

\subsubsection*{Obvyklý postup řešení}
Typické úlohy lze řešit podle následujícího schématu:

\begin{enumerate}[label=\arabic*.]
  \item Vyjádříme rovnici v obecnějším tvaru a najdeme kořen (nebo kořeny) v závislosti na parametru.
  \item Ošetříme \textbf{podmínky smysluplnosti} – například nesmí nastat dělení nulou, tedy určité hodnoty parametru je nutno vyloučit.
  \item Získáme kořen (nebo kořeny) jako funkci parametru.
  \item Na závěr uplatníme dodatečnou podmínku:  
    – kořen má být kladný,  
    – kořen má být $\ge$ určitá hodnota,  
    – kořen má ležet v daném intervalu,  
    – rovnice má mít právě jeden kořen (tedy \(D(p)=0\)).
\end{enumerate}

\subsubsection*{Ukázkový příklad}

\begin{examplebox}
\textbf{Najděte všechny hodnoty parametru \(a\), pro které má rovnice}
\[
4 - a = \frac{2}{x - 1}
\]
\textbf{kladný kořen.}

Nejprve upravíme rovnici na tvar bez zlomku:
\[
(4 - a)(x - 1) = 2.
\]

Za předpokladu, že \(4 - a \neq 0\), vydělíme:
\[
x - 1 = \frac{2}{4 - a},
\qquad
x = 1 + \frac{2}{4 - a}.
\]

Toto je kořen rovnice vyjádřený pomocí parametru.  
Aby byl kladný, musí platit
\[
1 + \frac{2}{4 - a} > 0.
\]

Řešením této nerovnice získáme přípustné intervaly hodnot parametru \(a\).  
Navíc musíme vyloučit hodnotu \(a = 4\), protože by došlo k dělení nulou.

\textbf{Výsledkem je množina všech hodnot parametru, pro které vyjde kladný kořen.}
\end{examplebox}

\begin{examplebox}
\textbf{Najděte všechna reálná čísla \(m\), pro která má rovnice}
\[
\frac{3x - 1}{m} - 5 = 2x
\]
\textbf{kořen splňující \(x \ge 4\).}

Rovnice je ekvivalentní tvaru
\[
3x - 1 = m(2x + 5).
\]

Za předpokladu \(m \neq 0\) vyjádříme kořen:
\[
3x - 1 = 2mx + 5m,
\]
\[
x(3 - 2m) = 5m + 1,
\]
\[
x = \frac{5m + 1}{3 - 2m},
\qquad
3 - 2m \neq 0.
\]

A protože má platit \(x \ge 4\), získáváme nerovnici
\[
\frac{5m + 1}{3 - 2m} \ge 4,
\]
kterou vyřešíme standardní metodou racionálních nerovnic.  
Součástí řešení je opět kontrola bodů, v nichž by docházelo k dělení nulou.
\end{examplebox}





\section{Kvadratické nerovnice}

Zatímco u kvadratických rovnic hledáme hodnoty, pro které se výraz rovná nule,  
u kvadratických nerovnic nás zajímá, kdy je kvadratický výraz \textit{větší} nebo \textit{menší} než nula.  
Jinými slovy, zkoumáme, v jakých intervalech je parabola nad osou \(x\) (kladné hodnoty) nebo pod osou \(x\) (záporné hodnoty).

Řešení kvadratické nerovnice spočívá v určení \textbf{nulových bodů} odpovídající rovnice \(ax^2 + bx + c = 0\)  
a následném zjištění, jaké znaménko má výraz \(ax^2 + bx + c\) v intervalech vymezených těmito body.  
Výsledkem je množina všech hodnot \(x\), pro které je nerovnost splněna.

\begin{definitionbox}
\textbf{Obecný tvar kvadratické nerovnice:}
\[
ax^2 + bx + c > 0, \quad ax^2 + bx + c \ge 0, \quad ax^2 + bx + c < 0, \quad ax^2 + bx + c \le 0,
\]
kde \(a, b, c \in \mathbb{R}\) a \(a \neq 0\).

Řešením je množina všech reálných čísel \(x\), pro která je výraz \(ax^2 + bx + c\) kladný, záporný nebo nulový,  
nejčastěji vyjádřená jako interval nebo sjednocení intervalů na číselné ose.
\end{definitionbox}

\begin{examplebox}
\textbf{Příklad:}  
Vyřešte kvadratickou nerovnici
\[
x^2 - 3x - 10 \le 0.
\]

Nejprve najdeme nulové body odpovídající rovnice
\[
x^2 - 3x - 10 = 0.
\]

Tento výraz lze rozložit na součin:
\[
(x - 5)(x + 2) = 0.
\]

Z toho vyplývá, že nulové body jsou
\[
x_1 = -2, \quad x_2 = 5.
\]

Nerovnici tedy můžeme přepsat jako
\[
(x - 5)(x + 2) \le 0.
\]

Součin dvou výrazů je záporný tehdy, když mají různé znaménko.  
Znamení jednotlivých faktorů určíme podle nulových bodů:

\[
\begin{array}{c|ccc}
x & (-\infty,-2) & (-2,5) & (5,\infty) \\ \hline
x+2 & - & + & + \\
x-5 & - & - & + \\ \hline
(x+2)(x-5) & + & - & +
\end{array}
\]

Protože hledáme, kdy je výraz \(\le 0\), tedy záporný nebo nulový,  
řešení tvoří interval mezi nulovými body včetně jejich konců.

\[
\textbf{Řešení:} \quad K = \langle -2;\,5 \rangle.
\]
\end{examplebox}




% -----------------------------------------------------------------
% ---------------------- P Ř Í K L A D Y --------------------------
% -----------------------------------------------------------------

\exercisepart
\setcounter{section}{1}

% =========================
% ÚLOHA 1 – Kvadratické rovnice k procvičení
% =========================
\begin{exercise}[Kvadratické rovnice – procvičování]

Vyřešte následující kvadratické rovnice v oboru reálných čísel \(x \in \mathbb{R}\).

\begin{tasks}(1)
    \task \(4x^2 + x = 0\)
    \task \(2x^2 - 5 = 0\)
    \task \(3x^2 + x - 2 = 0\)
    \task \(x^2 - 2x - 1 = 0\)
    \task \(4x^2 - x + 1 = 0\)
    \task \(3x^2 + x + 2 = 0\)
    \task \(x^2 + (2\sqrt{3} + 1)x + (3 + \sqrt{3}) = 0\)
    \task \(x^2 + x(2\sqrt{3} + 1) + 2\sqrt{3} = 0\)
    \task \(x^2 - \sqrt{2}\,x + x - \sqrt{2} = 0\)
    \task 
    $
    \frac12(2x - 1)^2 - \left[\frac12(x+1)\right]^2
    =
    3\left[\left(\frac12 x\right)^2 - \left(\frac12\right)^2\right]
    $
\end{tasks}
\end{exercise}


% =========================
% ÚLOHA – Ověření kořenů
% =========================
\begin{exercise}[Ověření, že zadané číslo je kořenem rovnice]
Ověřte, že zadané číslo je kořenem dané kvadratické rovnice.

\begin{tasks}(1)

  \task Ověřte, že číslo \(x = 2 + 3\sqrt{2}\) je kořenem rovnice  
  \[
  x^2 - 6x - 10 + 6\sqrt{2} = 0.
  \]

  \task Ověřte, že číslo \(x = 1 + 2\sqrt{3}\) je kořenem rovnice  
  \[
  x^2 - 2x - 11 - 4\sqrt{3} = 0.
  \]

  \task Ověřte, že číslo \(x = 4 - \sqrt{5}\) je kořenem rovnice  
  \[
  x^2 - 8x + 11 - 4\sqrt{5} = 0.
  \]

\end{tasks}

\end{exercise}


% =========================
% ÚLOHA – Sestavení kvadratické rovnice z kořenů
% =========================
\begin{exercise}[Sestavení kvadratické rovnice z kořenů]
Sestavte všechny kvadratické rovnice, jejichž kořeny jsou daná čísla.  
(Využijte vztah \(a(x - x_1)(x - x_2) = 0\), kde \(a \in \mathbb{R}, a \neq 0\).)

\begin{tasks}(2)
  \task \(x_1 = 2;\; x_2 = -3\)
  \task \(x_1 = -1;\; x_2 = -1\)
  \task \(x_1 = 1+\sqrt{3};\; x_2 = 1-\sqrt{3}\)
  \task \(x_1 = 0;\; x_2 = 5\)
  \task \(x_1 = \dfrac{1}{10-\sqrt{50}};\; x_2 = \dfrac{1}{10+5\sqrt{2}}\)
\end{tasks}

\end{exercise}


% =========================
% ÚLOHA – Určení koeficientů z daného kořene
% =========================
\begin{exercise}[Určení chybějícího koeficientu a druhého kořene]
V následujících rovnicích je jeden z kořenů známý. Určete druhý kořen a dopočítejte chybějící koeficient.

\begin{tasks}(2)

  \task Rovnice \(x^2 + px + 6 = 0\) má jeden kořen \(-2\).  
        Vypočítejte druhý kořen a koeficient \(p\).

  \task Rovnice \(x^2 - 3x + q = 0\) má jeden kořen \(-1\).  
        Vypočítejte druhý kořen a absolutní člen \(q\).

  \task Rovnice \(x^2 + rx - 12 = 0\) má jeden kořen \(4\).  
        Vypočítejte druhý kořen a koeficient \(r\).

  \task Rovnice \(x^2 + kx + 2 = 0\) má jeden kořen \(\tfrac{1}{2}\).  
        Vypočítejte druhý kořen a koeficient \(k\).

\end{tasks}

\end{exercise}


% =========================
% ÚLOHA – Parametrické kvadratické rovnice
% =========================
\begin{exercise}[Parametrické kvadratické rovnice]
Určete parametry tak, aby kořeny daných kvadratických rovnic splňovaly uvedené podmínky.

\begin{tasks}(1)

  \task Určete všechny hodnoty parametru \(b \in \mathbb{R}\) tak, aby jeden kořen kvadratické rovnice
  \[
  2x^2 + bx + 9 = 0
  \]
  byl dvakrát větší než druhý kořen.

  \task Určete koeficient lineárního členu \(m \in \mathbb{R}\) tak, aby jeden kořen kvadratické rovnice
  \[
  x^2 + mx + 5 = 0
  \]
  byl o \(4\) větší než druhý kořen.

  \task Určete všechny hodnoty absolutního členu \(q \in \mathbb{R}\) tak, aby jeden kořen kvadratické rovnice
  \[
  4x^2 - 15x + q = 0
  \]
  byl druhou mocninou druhého kořene.

\end{tasks}

\end{exercise}


% =========================
% ÚLOHA – Bonusová úloha se substitucí
% =========================
\begin{exercise}[Bonus: Řešení pomocí substituce]
Řešte v oboru reálných čísel rovnici
\[
3x^2 + 15x + 2\sqrt{x^2 + 5x + 1} = 2.
\]

Naznačte postup řešení pomocí vhodné substituce.

\end{exercise}


% =========================
% ÚLOHA – Rozklad kvadratického trojčlenu
% =========================
\begin{exercise}[Rozklad kvadratického trojčlenu na dvojčleny]
Rozložte následující kvadratické trojčleny na součin dvou lineárních výrazů (pokud je to možné).

\begin{tasks}(2)

  \task \(x^2 - 5x + 6\)
  \task \(x^2 - 7x + 10\)
  \task \(x^2 + x - 2\)
  \task \(x^2 - 11x + 30\)
  \task \(x^2 + 10x + 25\)
  \task \(4x^2 + 4x + 1\)
  \task \(16x^2 - 25\)
  \task \(15x^2 - 30x\)
  \task \(2x^2 + x - 1\)
  \task \(-x^2 + 7x + 8\)
  \task \(-2x^2 + 11x - 15\)
  \task \(x^2 + 4x + 5\)

\end{tasks}

\end{exercise}


% =========================
% ÚLOHA – Racionální rovnice
% =========================
\begin{exercise}[Racionální rovnice]
Řešte následující rovnice s neznámou \(x \in \mathbb{R}\).

\begin{tasks}(1)

  \task 
  \[
  \frac{3}{x+2} + \frac{5x}{4 - x^2} = \frac{3}{x-2} + \frac{x}{x^2 - 4}
  \]

  \task
  \[
  1 + 3\left( \frac12 \cdot \frac{x-3}{x-2} - \frac{2}{x-2} \right) = \frac{15}{2x - x^2}
  \]

  \task
  \[
  \left( \frac{x}{x-1} - \frac{x+1}{x} \right) : \left( \frac{x}{x+1} - \frac{x-1}{x} \right) = \frac{x+1}{x-1}
  \]

\end{tasks}

\end{exercise}


% =========================
% ÚLOHA – Rovnice s odmocninami
% =========================
\begin{exercise}[Rovnice s odmocninami]
Řešte následující rovnice v oboru reálných čísel \(x \in \mathbb{R}\).

\begin{tasks}(2)

  \task \(\sqrt{x} + x = 2\)

  \task \(3\sqrt{x+5} - 5 = x\)

  \task \(\sqrt{x-1} + \sqrt{x+4} = 5\)

  \task \(\sqrt{x+2} - 2\sqrt{x+7} = -4\)

  \task \(\sqrt{-x} = 2 - \sqrt{2 - x}\)

  \task \(\sqrt{10 - x} + \sqrt{x - 10} = 0\)

  \task \(\sqrt{2x + 7} + \sqrt{x - 5} = \sqrt{3x + 2}\)

  \task \(\sqrt{-x - \sqrt{1 - x}} = 1\)

  \task $
  \frac{2x}{2\sqrt{3} + \sqrt{x}} = 
  \frac{2\sqrt{3} - \sqrt{x}}{4 \frac{1}{2}}
  $

  \task 
  $
  \frac{4}{x + \sqrt{x^2 + x}} 
  - 
  \frac{1}{x - \sqrt{x^2 + x}} 
  = 
  \frac{3}{x}
  $

\end{tasks}

\end{exercise}


% =========================
% ÚLOHA – Rovnice řešené metodou substituce
% =========================
\begin{exercise}[Rovnice řešené metodou substituce]
Řešte následující rovnice v oboru reálných čísel \(x \in \mathbb{R}\) pomocí metody substituce.

\begin{tasks}(2)

  \task \(x^4 - 5x^2 + 4 = 0\)

  \task \(x^6 - 7x^3 - 8 = 0\)

  \task \(20x^{-4} + 3x^{-2} - 2 = 0\)

  \task \(5\sqrt[4]{x} - \sqrt{x} - 6 = 0\)

\end{tasks}

\end{exercise}


% =========================
% ÚLOHA – Rovnice vhodné pro substituci
% =========================
\begin{exercise}[Rovnice vhodné pro substituci]
Řešte následující rovnice v oboru reálných čísel \(x \in \mathbb{R}\).

\begin{tasks}(2)
  \task $ (1 - x^2)^2 - 2(x^2 - 1) + 1 = 0$
  \task $(x^2 + 2x - 3)(x^2 + 2x + 1) - 5 = 0$
  \task $(x^2 - x + 2)^2 - 6\bigl((x^2 - x) - 4\bigr) = 0$
  \task $(x^2 - 2x + 8)^2 : (x^2 - 2x)^2 - 5 = 0$
\end{tasks}

\end{exercise}


% =========================
% BONUS 1
% =========================
\begin{exercise}[BONUS: Optimalizace s nerovnostmi]
Kladná reálná čísla \(x, y, z\) splňují nerovnosti
\[
xy \ge 2,\qquad xz \ge 3,\qquad yz \ge 6.
\]
Jaké nejmenší hodnoty může nabývat výraz
\[
13x^2 + 10y^2 + 5z^2 \; ?
\]
\end{exercise}

% =========================
% BONUS 2
% =========================
\begin{exercise}[BONUS: Rovnost dvou kvadratických výrazů]
Předpokládejme, že pro reálná čísla \(a, b\) mají výrazy
\[
a^2 + b \qquad \text{a} \qquad a + b^2
\]
stejnou hodnotu.
Jaká nejmenší může být tato společná hodnota?
\end{exercise}

% =========================
% BONUS 3
% =========================
\begin{exercise}[BONUS: Soustava kvadratických rovnic]
V oboru reálných čísel řešte soustavu
\[
\begin{cases}
xy + 1 = z^2,\\
yz + 2 = x^2,\\
zx + 3 = y^2.
\end{cases}
\]
\end{exercise}


% =========================
% ÚLOHA – Pohyb pístu (poloha)
% =========================
\begin{exercise}[Pohyb pístu – poloha jako kvadratická funkce]
V ideálním modelu se poloha pístu ve válci (v metrech) popisuje v závislosti na čase \(t\) (v sekundách) funkcí
\[
s(t) = -0{,}5t^2 + 0{,}6t,
\]
kde \(t \in \langle 0; 1{,}2 \rangle\). Hodnota \(s(t)\) udává vzdálenost pístu od výchozí polohy.

\begin{enumerate}[label=\alph*)]
  \item Určete souřadnice vrcholu této paraboly a vysvětlete, co fyzikálně znamenají.
  \item V jakém čase je píst nejdále od výchozí polohy a jaká je tato maximální vzdálenost?
  \item Pro jaké hodnoty času je píst dále než \(0{,}25\) m od výchozí polohy?
  \item V jakých okamžicích se píst nachází ve výchozí poloze? Co to znamená z hlediska jednoho pracovního zdvihu?
\end{enumerate}

\end{exercise}

% =========================
% ÚLOHA – Pohyb pístu (rychlost)
% =========================
\begin{exercise}[Pohyb pístu – rychlost jako kvadratická funkce]
Rychlost pístu ve válci (v metrech za sekundu) popisujeme v čase \(t\) (v sekundách) funkcí
\[
v(t) = -80t^2 + 64t,
\]
kde \(t \in \langle 0; 0{,}8 \rangle\).

\begin{enumerate}[label=\alph*)]
  \item Určete souřadnice vrcholu grafu funkce \(v(t)\) a vysvětlete, co fyzikálně znamenají.
  \item V jakých okamžicích je rychlost pístu nulová a jaký význam mají tyto body pro pohyb pístu?
  \item V jakém intervalu časů je rychlost pístu kladná a v jakém záporná? Co to znamená pro směr pohybu?
  \item Určete maximální rychlost pístu během sledovaného intervalu.
\end{enumerate}

\end{exercise}

\setcounter{section}{2}
% =========================
% ÚLOHA – Rozměry garáže
% =========================
\begin{exercise}[Rozměry obdélníkové garáže]
Obdélníková garáž má plochu \(24\,\text{m}^2\).  
Jedna její strana je o \(2\ \text{m}\) delší než druhá.  
Určete rozměry garáže.

\end{exercise}


% =========================
% ÚLOHA – Rozměry obdélníkového pozemku
% =========================
\begin{exercise}[Rozměry obdélníkového pozemku]
Obdélníkový pozemek má jednu stranu o polovinu delší než druhou.  
Určete jeho rozměry, jestliže má plochu \(96\,\text{a}\).

\end{exercise}


% =========================
% ÚLOHA – Rozměry lichoběžníku
% =========================
\begin{exercise}[Rozměry lichoběžníku]
Jedna ze základen lichoběžníku je o pětinu větší než jeho výška,  
druhá základna je o \(1\ \text{cm}\) větší než jeho výška.  
Určete rozměry tohoto lichoběžníku, jestliže jeho obsah je
$2115\ \text{cm}^2.$

\end{exercise}


% =========================
% ÚLOHA – Poloměr kruhu před zvětšením
% =========================
\begin{exercise}[Poloměr kruhu před zvětšením]
Když jsme poloměr kruhu zvětšili o \(2\ \text{cm}\), zvětšil se jeho obsah o
$
240\pi\ \text{cm}^2.
$
Určete poloměr kruhu před zvětšením.

\end{exercise}


% =========================
% ÚLOHA – Plnění vodní nádrže
% =========================
\begin{exercise}[Plnění vodní nádrže třemi přívody]
Vodní nádrž je možno plnit třemi přívody. Druhým přívodem by se nádrž naplnila
o polovinu delší dobu než prvním, třetím přívodem pak o \(8\) hodin déle než prvním.
Všemi třemi přívody současně se nádrž naplní za \(4{,}5\) hodiny.

Za jakou dobu by se nádrž naplnila jednotlivými přívody?

\end{exercise}


% =========================
% ÚLOHA – Rychlost vlaku
% =========================
\begin{exercise}[Rychlost vlaku a doba jízdy]
Pokud se rychlost vlaku zvýší o \(9\ \text{km·h}^{-1}\), urazí dráhu \(180\ \text{km}\)
o \(40\) minut dříve než při původní rychlosti.

Vypočítejte čas, za který by vlak ujel tuto dráhu při původní rychlosti.

\end{exercise}


% =========================
% ÚLOHA – Dvojciferné číslo
% =========================
\begin{exercise}[Dvojciferné číslo se zadanými vlastnostmi]
Najděte dvojciferné číslo, pro které platí:

\begin{itemize}
  \item ciferný součet je \(9\),
  \item po prohození číslic vznikne nové číslo, které po vynásobení původním číslem
        dává součin \(2430\).
\end{itemize}

\end{exercise}


% =========================
% ÚLOHA – Dva rezistory v obvodu
% =========================
\begin{exercise}[Určení odporů dvou rezistorů]
V obvodu, v němž jsou zapojeny paralelně dva rezistory, prochází při napětí 
\(24\ \text{V}\) proud \(4\ \text{A}\).  
Pokud tyto rezistory zapojíme sériově, klesne proud na \(0{,}75\ \text{A}\).

Určete odpory obou rezistorů.

\end{exercise}


% =========================
% ÚLOHA – Kvadratické nerovnice
% =========================
\begin{exercise}[Řešte kvadratické nerovnice]
Řešte následující kvadratické nerovnice v oboru reálných čísel.

\begin{tasks}(2)

  \task \(x^2 - 1 \ge 0\)
  \task \(x^2 - 16 \le 0\)
  \task \(-x^2 - 4x - 3 \ge 0\)
  \task \(x^2 - 1 < 0\)

  \task \(x^2 - 9x + 14 \ge 0\)
  \task \(x^2 + 4x + 9 \le 0\)
  \task \(-x^2 - 2x + 8 > 0\)
  \task \(x^2 - 6x + 7 < 0\)

  \task \(x^2 - x - 12 < 0\)
  \task \(2x^2 - 3x + 1 > 0\)
  \task \(2x^2 + 3x + 4 > 0\)
  \task \(x^2 - 8x + 12 \ge 0\)

  \task \(4x^2 > 12x\)
  \task \(4x - x^2 > 12\)
  \task \(x(x - 2) \le 2x - 4\)
  \task \(x^2 + 2x - 3 \ge 0\)

  \task \(x^2 + 6 > -2x\)
  \task \((x - 1)(x - 2) > 0\)
  \task \((x + 3)(x - 5) \le 0\)
  \task \(2x^2 - 5x - 7 > 0\)

  \task \(2x^2 - x - 6 < 0\)
  \task \(x^2 - x - 1 \ge 0\)
  \task \(x^2 + 2x + 1 \le 0\)
  \task \(x^2 + 8x + 16 \ge 0\)

  \task \(x^2 - 4x - 12 \ge 0\)
  \task \(-2(x^2 - 1) + 2(1 - x) - 5 \le 0\)

\end{tasks}

\end{exercise}


% =========================
% ÚLOHA – Racionální nerovnice
% =========================
\begin{exercise}[Řešte racionální nerovnice]
Řešte následující nerovnice v oboru reálných čísel.

\begin{tasks}(2)

\task \(\displaystyle \frac{x^4 + x^2 + 1}{x^2 - 4x - 5} < 0\)

\task \(\displaystyle \frac{x^2 + 4x + 4}{2x^2 - x - 1} > 0\)

\task \(\displaystyle \frac{x^2 - 5x + 6}{x^2 + x + 1} < 0\)

\task \(\displaystyle \frac{1 - 2x - 3x^2}{3x - x^2 - 5} > 0\)

\task \(\displaystyle \frac{x^2 - 5x + 12}{x^2 - 4x + 5} > 3\)

\task \(\displaystyle \frac{x^2 - 3x + 24}{x^2 - 3x + 3} < 4\)

\task \(\displaystyle \frac{4}{1 + x} + \frac{2}{1 - x} < 1\)

\task \(\displaystyle \frac{x^4 - 3x^3 + 2x^2}{x^2 - x - 30} > 0\)

\task \(\displaystyle \frac{2(x - 3)}{(x - 1)(x - 10)} \ge \frac{1}{x - 2}\)

\task \(\displaystyle \frac{7}{(x - 2)(x - 3)} + \frac{9}{x - 3} < -1\)

\task \(\displaystyle \frac{(x - 1)(x + 2)(2 - x)}{(3 - x)(x - 4)} \le 0\)

\task \(\displaystyle \frac{(x - 1)^2 \cdot 5x (x + 2)}{(x - 4)(x - 5)^2} \ge 0\)

\end{tasks}

\end{exercise}



% =========================
% ÚLOHA 4
% =========================
\begin{exercise}[Parametrická rovnice – kladný kořen]
Najděte všechna reálná čísla \(a\) tak, aby daná rovnice měla kladné řešení:
\begin{tasks}
  \task $ 4 - a = \dfrac{2}{x - 1}$
  \task $ x \geq 4:\dfrac{3x - 1}{m} - 5 = 2x$
  \task $ \dfrac{5x - 3}{p} + 2 = x$
  \task $ 4x - k = \dfrac{8}{x - 1}$
\end{tasks}
\end{exercise}



\end{document}